\sscp{Makuva}{04-06-01}
\ili{Makuva} [lva] (a.k.a \ili{Lovaia}) is a highly endangered language of
eastern Timor which is reportedly only spoken by a few families
in the Lautém district on the far eastern tip of Timor.
\citet{vanEngelenhoven2009} provides a summary of the top-down
historical phonology of \ili{Makuva},
which attests several unusual sound changes such as *t > \it{k}
(also attested in other \tsc{Southwest Maluku} languages)
and *s > \it{t}.
It also has several conservative features such as retention of
PMP *z [ʤ] = \it{ʤ}, initial *d- = \it{d}, and *R = \it{r}.
Examples of the principle sound correspondences between PMP
and \ili{Makuva} are given in \trf{04-26}.

%\begin{table}\tcp{PMP to \ili{Makuva} correspondences (van Engelenhoven 2009:431--434)}{04-26}
%\begin{table}\caption{PMP to Makuva correspondences \citep[431ff]{vanEngelenhoven2009}}\label{tab:04-26}
\begin{table}[p]\centering\caption{Makuva correspondences \citep[431ff]{vanEngelenhoven2009}}\label{tab:04-26}
\begin{adjustbox}{totalheight=\textheight}
	\st{3pt}\wg\begin{tabularx}{\linewidth}{llllX}\lsptoprule
\mc{2}{l}{PMP}		&						&	env.				&	Examples	\\	\midrule
*p	&	>	&	\it{h}		&	{\#}{\gap}	&	*paŋdan > \it{hene} ‘pandanus’, *pusəj > \it{hutre} ‘navel’ 	\\	
		&	>	&	{\0}			&	/V{\gap}V		&	*nipay > \it{ne} ‘snake’, *kapuR > \it{βaro, waru} ‘calcium’	\\
*t	&	>	&	\it{k}		&							&	*təlu > \it{-kelu} ‘three’, *um-utaq > \it{muke} ‘vomit’, *kulit > \it{ulke} ‘skin’	\\	
*k	&	>	&	{\0}			&							&	*kutu > \it{uku} ‘louse’, *kahiw > \it{ai} ‘wood’, *ikan > \it{yene} ‘fish’	\\	
*q	&	>	&	{\0}			&							&	*taqi > \it{ke} ‘dung’, *qənay > \it{ini} ‘sand’, *qatay > \it{ake} ‘liver’	\\	
*b	&	>	&	\it{h}		&	{\#}{\gap}	&	*batu > \it{haku} ‘stone’, *bəli > \it{heli} ‘buy’	\\	
		&	>	&	{\0}			&	/V{\gap}V		&	*təbuh > \it{kiu} ‘sugarcane’, *qabu > \it{au} ‘ash’	\\	
*d	&	>	&	\it{d, t}	&	{\#}{\gap}	&	*daRaq > \it{dar-βe} ‘blood’, *dəŋəR > \it{dete {\tl} tete} ‘hear’	\\	
		&	>	&	\it{r}		&	/V{\gap}V, /{\gap}{\#}	&	*si-ida > \it{tira} ‘\tsc{3pl}’, *likud > \it{liru} ‘back’	\\	
*z	&	=	&	\it{ʤ}		&		&	*zalan > \it{ʤane}, *zaqat > \it{ʤake} ‘bad’, *quzan > \it{ʤone} ‘rain’	\\	
*j	&	>	&	\it{d, r}	&		&	*ŋajan > \it{nade-} ‘name’, *pusəj > \it{hutre,} ‘navel’, *laləj > \it{lare} ‘fly’	\\	
*ñ	&	>	&	\it{n}		&		&	*miñak > \it{mine} ‘fat’	\\	
*ŋ	&	>	&	\it{n}		&		&	*ŋisi > \it{nit-βe} ‘tooth’, *haŋin > \it{ane} ‘wind’, *taliŋa > \it{linu} ‘ear’	\\	
*s	&	>	&	\it{t}		&		&	*susu > \it{tutu} ‘br\ili{east}’, *sakay > \it{tai} ‘go up’, *asu > \it{ato} ‘dog’	\\	
*h	&	>	&	{\0}			&		&	*haŋin > \it{ane} ‘wind’, *wahiyəR > \it{were} ‘water’	\\	
*l	&	=	&	\it{l}		&		&	*laRiw > \it{lari, lori}, *balik > \it{hali} ‘return’, *qulu > \it{ulu} ‘head’	\\	
*R	&	=	&	\it{r}		&		&	*bəRəqat > \it{herka} ‘heavy’, *uRat > \it{urke} ‘vein’, *ikuR > \it{iru} ‘tail’	\\	
*i	&	=	&	\it{i}		&	/{\#}(C){\gap}	&	*ina > \it{in-} ‘mother’, *isa > \it{ite-tla} ‘one’	\\	
		&	>	&	\it{e}		&	/{\gap}C{\#}		&	*diŋdiŋ > \it{rine} ‘cold’, *kulit > \it{ulke} ‘skin’	\\	
*ə	&	>	&	\it{e}		&									&	*qaləjaw > \it{lera, lere}, *ma-qitəm > \it{mekme}	\\	
*a	&	=	&	\it{a}		&	/{\#}(C){\gap}	&	*tasik > \it{kate} ‘sea’, *anak > \it{ana, ane} ‘child’	\\	
		&	>	&	\it{a, u}	&	/{\gap}C{\#}		&	*mata > \it{maku, maka} ‘eye’, *lima > \it{-lima} ‘five’, *taliŋa > \it{linu} ‘ear’	\\	
*u	&	=	&	\it{u}		&		&	*duha > \it{-rua} ‘two’, *qulu > \it{ulu} ‘head’	\\	
*-iw&	>	&	\it{i}		&		&	*laRiw > \it{lari, lori} ‘run’, *kahiw > \it{ai} ‘wood’	\\	
*-ay&	>	&	\it{e, i}	&		&	*matay > \it{make} ‘dead’, *qatay > \it{ake} ‘liver’, *qənay > \it{ini} ‘sand’	\\	
*-aw&	>	&	\it{a, e}	&		&	*bətaw > \it{maheka, maheke} ‘female’, *qaləjaw > \it{lera, lere} ‘sun’	\\	
		\lspbottomrule
	\end{tabularx}
\end{adjustbox}
\end{table}

\ili{Makuva} shows extensive evidence of final VC > CV metathesis which
reflects the changes in \qf{04-12} (\prf {ex:04-12}) which define \tsc{Southwest Maluku}.
However, such metathesis is often blurred by subsequent simplification
of medial consonant clusters as well as apparent final *ə > \it{e, a}.
Examples in addition to those in \trf{04-25} include: PMP *pus\tbf{əj}
> \it{hut\tbf{re}} ‘navel’, *maqit\tbf{əm} > \it{mek\tbf{me}} ‘black’,
*kul\tbf{it} ‘skin’ > \it{ul\tbf{ke}},
and regional **mal\tbf{us} > \it{mal\tbf{tu}} ‘betel pepper’
(e.g. \ili{Tetun} \it{malus}, \tsc{Meto} \it{manus},
\ili{Kisar} \it{maluh/malhu}).

While some of the changes in \ili{Makuva} also take place in some other
languages of \tsc{Southwest Maluku} (e.g.
*t > \it{k} in \ili{Kisar},
*s > \it{t} in \tsc{South Babar}),
no other shared innovations unambiguously occur in both \ili{Makuva}
and any of the other primary nodes of \tsc{Southwest Maluku}.
Thus, it is placed as an isolate within \tsc{Southwest Maluku}.
